%
% Draft paper in EAMT-style LaTeX.
% If eamt26.sty is unavailable locally, this file can still be copied into
% the official EAMT template directory and compiled there.
%
\documentclass[11pt]{article}
\usepackage{eamt26}
\usepackage[utf8]{inputenc}
\usepackage[T1]{fontenc}
\usepackage{times}
\usepackage{url}
\usepackage{latexsym}
\usepackage[small,bf]{caption}
\usepackage{booktabs}
\usepackage{multirow}
\usepackage{array}
\usepackage{graphicx}
\usepackage{amsmath}
\setlength\titlebox{6.5cm}
\usepackage{tikz}
\usetikzlibrary{shapes,arrows.meta,positioning}
\usepackage{pgfplots}
\usepgfplotslibrary{groupplots}
\pgfplotsset{compat=1.18}
\usepackage[hidelinks]{hyperref}

\newcommand{\confname}{EAMT 2026}
\newcommand{\conffilename}{eamt26}
\newcommand{\nb}{Bokm{\aa}l}
\newcommand{\nn}{Nynorsk}

\title{When Data Cleaning Becomes Bias: Target-Standard Specialization in Norwegian Machine Translation}

\author{Anonymous submission}

\date{}

\begin{document}
\maketitle

\begin{abstract}
Languages with multiple written standards pose a challenge for machine translation (MT): preprocessing decisions that appear linguistically innocuous may systematically favor one legitimate standard over another. Norwegian is a clear example. Although \nb{} and \nn{} are both official written standards, MT pipelines and evaluation references often implicitly privilege one of them. We refer to this phenomenon as \emph{target-standard bias}: a systematic preference for one written standard that can be introduced during data preparation and amplified during evaluation.

We investigate target-side \nb{} filtering for English--Norwegian petroleum-domain MT. Using LoRA, we fine-tune NLLB-200 models on original mixed-standard training data, an \nb{}-filtered subset, and a size-controlled random subset. Across 600M, 1.3B, and 3.3B models, \nb{} filtering improves \nb{}-conforming automatic performance under \nb{} references relative to the size-controlled baseline, but reduces BLEU and chrF on the original mixed-standard test set.

SLIDE-based analysis shows the same scale-stable pattern: on the original mixed-standard test set, \nb{} filtering increases \nb{}-only outputs from about 79--80\% to 94--95\% across model sizes while nearly eliminating \nn{}-only outputs. Manual validation confirms the main \nb{}/\nn{} contrast but also shows that mixed or weakly diagnostic cases are genuinely ambiguous, so automatic written-standard labels should be treated as diagnostic evidence rather than a fine-grained gold standard. Glossary-based terminology metrics show a parallel increase in \nb{} term-form conformity rather than language-neutral terminology accuracy. Out-of-domain FLORES evaluation for the 600M systems indicates that these effects are largely domain-specific, and that low scores against \nn{} references in this setup mainly reflect written-standard mismatch rather than \nn{} translation ability.

Our results show that target-side filtering is not a neutral data-cleaning step. Instead, it induces target-standard specialization, and automatic evaluation metrics may reward this specialization when reference data favor the same written standard.
\end{abstract}

\begin{figure}[t]
\centering
\resizebox{\columnwidth}{!}{%
\begin{tikzpicture}[
  every node/.style={font=\small},
  corpusbox/.style={rectangle, draw, rounded corners, align=center, minimum height=0.8cm, fill=gray!15, inner sep=5pt},
  databox/.style={rectangle, draw, rounded corners, align=center, minimum height=0.7cm, minimum width=2.6cm, fill=blue!10, inner sep=4pt},
  trainbox/.style={rectangle, draw, rounded corners, align=center, minimum height=0.8cm, fill=gray!15, inner sep=5pt},
  evalbox/.style={rectangle, draw, rounded corners, align=center, minimum height=0.7cm, minimum width=2.6cm, fill=green!10, inner sep=4pt},
  rqbox/.style={rectangle, draw, dashed, align=center, minimum height=0.6cm, minimum width=2.6cm, fill=yellow!12, inner sep=3pt},
  arr/.style={-{Latex[length=1.5mm]}}
]

% Row 1: corpus
\node[corpusbox, minimum width=8cm] (corpus) at (0,0) {Original EN--NO Corpus (Mixed \nb{}/\nn{})};

% Row 2: three data conditions
\node[databox] (orig) at (-2.8,-1.5) {Original Data};
\node[databox] (nb)   at (0,-1.5)    {\nb{}-Filtered\\(SLIDE)};
\node[databox] (rand) at (2.8,-1.5)  {Original-Sub.\\(size-matched)};

% Row 3: fine-tuning
\node[trainbox, minimum width=8cm] (lora) at (0,-3.0) {LoRA Fine-Tuning\\(NLLB-200; 600M/1.3B/3.3B scaling check)};

% Row 4: three systems
\node[trainbox, minimum width=8cm] (systems) at (0,-4.5) {Three Fine-Tuned MT Systems};

% Row 5: evaluation stages
\node[evalbox] (indomain) at (-2.8,-6.0) {In-Domain\\BLEU/chrF/Term};
\node[evalbox] (slide)    at (0,-6.0)    {SLIDE\\Standard Analysis};
\node[evalbox] (flores)   at (2.8,-6.0)  {FLORES\\Out-of-Domain};

% Row 6: RQ labels
\node[rqbox] (rq123) at (-2.8,-7.2) {RQ1--3, RQ6};
\node[rqbox] (rq4)   at (0,-7.2)    {RQ4};
\node[rqbox] (rq5)   at (2.8,-7.2)  {RQ5};

% Arrows
\draw[arr] (corpus) -- (orig);
\draw[arr] (corpus) -- (nb);
\draw[arr] (corpus) -- (rand);

\draw[arr] (orig) -- (lora);
\draw[arr] (nb) -- (lora);
\draw[arr] (rand) -- (lora);

\draw[arr] (lora) -- (systems);

\draw[arr] (systems) -- (indomain);
\draw[arr] (systems) -- (slide);
\draw[arr] (systems) -- (flores);

\draw[arr] (indomain) -- (rq123);
\draw[arr] (slide) -- (rq4);
\draw[arr] (flores) -- (rq5);

\end{tikzpicture}%
}
\caption{Experimental pipeline. Dashed boxes indicate which research questions each evaluation stage addresses. RQ6 (model-scale replication) uses the same in-domain automatic-metric setup as RQ1--RQ3.}
\label{fig:pipeline}
\end{figure}

\section{Introduction}

Machine translation systems are usually evaluated as if each target language corresponds to a single written norm. This assumption is convenient but not always true. Norwegian has two official written standards, \nb{} and \nn{}, which differ in spelling, morphology, lexical choice, and sociolinguistic distribution. A translation system configured to generate Norwegian may therefore still need to decide which written standard to produce. This decision can be explicit, as in a target language code such as \texttt{nob\_Latn}, or implicit, as in a training corpus whose target side contains mostly one standard.

This paper studies a practical data-curation decision: filtering a domain-specific English--Norwegian corpus toward \nb{}. At first glance, this looks like ordinary cleaning. However, filtering also changes the target distribution. We define \emph{target-standard bias} as a systematic preference for one legitimate written standard over another, introduced through data preparation and reinforced by evaluation. For multi-standard languages, target-side filtering may therefore influence both model behavior and evaluation outcomes. Section~\ref{sec:bias} formalizes this framing before the empirical analysis.

We investigate this question in English--Norwegian petroleum-domain MT. Domain MT is a relevant setting because terminology consistency matters and specialized corpora are often curated before fine-tuning. We compare three systems: a model fine-tuned on the original data, a model fine-tuned on an \nb{}-filtered subset, and a model fine-tuned on a random original subset of the same size as the filtered data. The size-controlled system allows us to separate filtering effects from simple reductions in training data.

Our experiments examine how target-side filtering affects automatic MT scores, terminology behavior, written-standard distributions, and robustness across model sizes, with an additional out-of-domain FLORES check. By combining automatic metrics, terminology evaluation, written-standard identification, and diagnostic human assessment, we analyze both the performance consequences of filtering and the target-standard assumptions embedded in MT evaluation. Figure~\ref{fig:pipeline} summarizes how the data conditions, model-scale fine-tuning, and evaluation stages connect to these questions.

The paper contributes:
\begin{enumerate}
    \item We introduce \emph{target-standard bias} as an evaluation problem for MT in multi-standard languages.
    \item We show that \nb{} filtering induces dominant-standard specialization beyond data-size effects across multiple NLLB-200 model sizes.
    \item We demonstrate that written-standard mismatch can reverse or complicate the interpretation of BLEU, chrF, and terminology metrics.
    \item We combine automatic written-standard identification with human evaluation to distinguish translation adequacy from standard conformity.
    \item We discuss implications for linguistic inclusion, deployment transparency, and MT evaluation in multi-standard language settings.
\end{enumerate}

\section{Target-Standard Bias}
\label{sec:bias}

In this paper, bias is used in the evaluation sense rather than to denote demographic or political stereotypes. We define \emph{target-standard bias} as a systematic preference for one legitimate written standard over another, introduced through data curation, model configuration, or evaluation design.

For Norwegian MT, target-standard bias arises when a system described simply as ``Norwegian'' systematically favors one official written standard, typically \nb{}, while valid \nn{} alternatives receive less support or lower evaluation scores.

Such bias can enter the MT pipeline at multiple stages. Target-side filtering may alter the distribution observed during training by removing minority-standard examples. Language codes such as \texttt{nob\_Latn} may encourage a particular written standard during decoding. Reference translations and terminology glossaries may further reinforce the same preference during evaluation by rewarding one standard while penalizing equally valid alternatives.

These choices can lead to \emph{silent specialization}: a system becomes increasingly aligned with a dominant written standard while continuing to be presented as general Norwegian MT. Because automatic metrics are computed against fixed reference forms, such specialization may appear as a quality improvement rather than a shift in written-standard preference.

\section{Research Questions}

To investigate target-standard bias empirically, we ask:

\begin{description}
\item[RQ1] Does \nb{} filtering improve performance on in-domain \nb{} test data?

\item[RQ2] Are the observed effects caused by filtering itself rather than the smaller size of the filtered dataset?

\item[RQ3] Does \nb{} filtering reduce robustness to the original mixed-standard distribution?

\item[RQ4] Does filtering change the written-standard distribution of generated outputs?

\item[RQ5] Do in-domain gains transfer to general-domain FLORES evaluation?

\item[RQ6] Does the size-controlled evaluation pattern persist across larger NLLB-200 model sizes?
\end{description}



\section{Related Work}

\paragraph{Language varieties and multi-standard MT.}
MT systems often under-serve language varieties that are not dominant in training data. Work on standard-to-variety or variety-aware MT has shown that closely related standards cannot always be treated as interchangeable, even when they are mutually intelligible. Kumar et al.~\shortcite{kumar2021machine}, for example, study adaptation from standard MT systems to related target varieties and include English--\nn{} as a case derived from English--\nb{}. Our work asks a complementary question: what happens when a domain-specific English--Norwegian MT pipeline is curated toward \nb{} rather than adapted toward \nn{}?

\paragraph{Variant and standard bias beyond Norwegian.}
The Norwegian case is part of a broader family of MT problems involving national, regional, or written-standard varieties. Lakew et al.~\shortcite{lakew-etal-2018-neural} show that treating varieties as a single target class can produce inconsistent translations, studying European/Brazilian Portuguese, European/Canadian French, Croatian/Serbian, and Indonesian/Malay. Costa-juss{\`a} et al.~\shortcite{costa-jussa-etal-2018-neural} focus specifically on translation between Brazilian and European Portuguese, and recent European Portuguese MT work similarly motivates variety-specific models by noting the dominance of Brazilian Portuguese resources \cite{sousa-etal-2025-tradutor}. Arabic MT provides another example where Modern Standard Arabic and regional dialects require explicit task and evaluation design \cite{nacar-etal-2024-asos,mousi-etal-2025-aradice}. Chinese localization tools such as OpenCC distinguish simplified/traditional scripts and regional lexical conventions for Mainland China, Taiwan, and Hong Kong \cite{opencc}. These cases differ linguistically from Norwegian, but they support the same methodological point: MT systems and evaluations should make the intended target variety or standard explicit.

\paragraph{Norwegian NLP.}
The \nb{}/\nn{} distinction is a central issue in Norwegian NLP. Norwegian language-resource and language-modeling work explicitly evaluates resources or models for both \nb{} and \nn{} \cite{kummervold2021large,kummervold-etal-2022-norwegian,samuel-etal-2025-small}. This is not only a linguistic matter but also a resource issue: \nb{} is more prevalent in many digital corpora, while \nn{} is lower-resource. Our study focuses on MT evaluation and domain adaptation rather than language model pretraining, but the same issue appears: target-side standard choice affects what the model learns to output.

\paragraph{Data curation as an explicit modeling choice.}
Bender and Friedman~\shortcite{bender-friedman-2018-data} argue that dataset documentation should expose the assumptions encoded in data collection and preprocessing, partly because such assumptions can contribute to system bias. We follow this perspective. Filtering toward \nb{} should not be reported as a minor preprocessing detail; it defines the target norm of the system. Our experiments measure the consequences of that choice as a form of target-standard bias.

\paragraph{Models, adaptation, and evaluation.}
We fine-tune NLLB-200 \cite{nllbteam2024scaling} with LoRA \cite{hu2021lora}. For general-domain out-of-domain evaluation, we use FLORES-200 \cite{nllbteam2024scaling}. For written-standard analysis, we use SLIDE \cite{fedorova2025slide}, which is designed for multi-label Scandinavian language identification and can mark sentences as compatible with more than one standard.

\section{Data}

\subsection{In-domain corpus}

The in-domain data consists of English--Norwegian sentence pairs from the petroleum domain. We use an original split and a \nb{}-filtered split. The original data contains a mixture of \nb{}, \nn{}, and mixed-standard Norwegian target sentences. The \nb{} split is obtained by applying SLIDE to the Norwegian target side and retaining examples with \nb{} score at least 0.80 and \nn{} score below 0.30. In the resulting \nb{} training split, all retained examples have higher \nb{} than \nn{} score; the minimum retained \nb{} score is 0.8021 and the mean is 0.9982. The resulting filtered data is smaller than the original data.

The filtering thresholds are intended to identify clear \nb{} examples rather than optimize evaluation outcomes. As a data-selection sensitivity check, nearby threshold settings retained between 9,925 and 10,212 training sentences, or 71.2--73.3\% of the original training split. This indicates that the chosen 0.80/0.30 threshold lies in a relatively stable region of the SLIDE score distribution. Details are reported in Appendix~\ref{app:threshold-sensitivity}. Because SLIDE is also used for diagnostic written-standard analysis, we validate its labels against human annotations in Section~\ref{sec:slide-analysis}.

This size difference is a methodological problem: if the \nb{} model behaves differently, the cause could be either the written-standard filtering or the change in training size. We therefore construct a size-controlled baseline, \emph{original-subsampled}, by randomly sampling the original train, validation, and test splits to match the \nb{} split sizes shown in Table~\ref{tab:data}. Sampling is deterministic with a fixed seed.

\begin{table}[t]
\centering
\small
\begin{tabular}{lrrr}
\toprule
Condition & Train & Valid. & Test \\
\midrule
Original & 13,935 & 1,737 & 1,742 \\
\nb{}-filtered & 10,114 & 1,305 & 1,313 \\
Original-subsampled & 10,114 & 1,305 & 1,313 \\
\bottomrule
\end{tabular}
\caption{In-domain dataset sizes. Original-subsampled controls for the smaller size of the \nb{}-filtered data.}
\label{tab:data}
\end{table}

\subsection{Out-of-Domain FLORES Data}

For out-of-domain evaluation, we use FLORES-200 devtest with
\texttt{eng\_Latn}$\rightarrow$\texttt{nob\_Latn}
(FLORES NB) and
\texttt{eng\_Latn}$\rightarrow$\texttt{nno\_Latn}
(FLORES NN) references. Each test set contains 1,012 examples.

FLORES differs substantially from the petroleum-domain corpus in topic, style, and vocabulary. It therefore provides a useful test of whether the effects observed in domain-specific fine-tuning transfer beyond the training domain.

FLORES NB is used to evaluate general-domain translation against \nb{} references, whereas FLORES NN provides a controlled written-standard mismatch condition by evaluating the same outputs against \nn{} references.

All systems in this study decode using the target language code \texttt{nob\_Latn}. Consequently, evaluation against FLORES NN should not be interpreted as a direct measure of \nn{} generation ability. Rather, it measures the evaluation penalty incurred when \nb{}-oriented outputs are compared against \nn{} references.

\section{Experimental Setup}

We fine-tune NLLB-200 models \cite{nllbteam2024scaling} with LoRA \cite{hu2021lora}. The main diagnostic experiments use \texttt{facebook/nllb-200-distilled-600M}. To test whether the observed size-controlled metric pattern is specific to this backbone, we also fine-tune \texttt{facebook/nllb-200-1.3B} and \texttt{facebook/nllb-200-3.3B} for the in-domain model-scale replication.

All fine-tuned systems use the same three data conditions: the full original corpus, the \nb{}-filtered corpus, and the size-matched original-subsampled corpus. We also evaluate the zero-shot NLLB base model without fine-tuning. Table~\ref{tab:model-conditions} summarizes the model conditions, and Table~\ref{tab:training-config} reports the shared fine-tuning configuration.

\begin{table}[!ht]
\centering
\small
\begin{tabular}{p{0.27\linewidth}p{0.63\linewidth}}
\toprule
Condition & Description \\
\midrule
Base & Zero-shot NLLB, no fine-tuning \\
Original & Fine-tuned on the full original training data \\
Original-sub. & Fine-tuned on the size-matched original subset \\
\nb{} & Fine-tuned on the \nb{}-filtered training data \\
\bottomrule
\end{tabular}
\caption{Model conditions evaluated in the study.}
\label{tab:model-conditions}
\end{table}

\begin{table}[t]
\centering
\small
\begin{tabular}{>{\raggedright\arraybackslash}p{0.30\linewidth} >{\raggedright\arraybackslash}p{0.60\linewidth}}
\toprule
Setting & Value \\
\midrule
LoRA rank / $\alpha$ / dropout & 8 / 64 / 0.0 \\
LoRA target modules & \texttt{q\_proj}, \texttt{k\_proj}, \texttt{v\_proj}, \texttt{out\_proj} \\
Epochs / learning rate & 3 / $5\times10^{-4}$ \\
Batching & batch size 4, gradient accumulation 4 \\
Warmup / checkpoint interval & 100 steps / 200 steps \\
Early stopping patience & 3 \\
Maximum sequence length & 128 \\
Training precision / optimizer & FP16 / AdamW, weight decay 0.01 \\
Generation & beam size 5, maximum length 128 \\
Seeds & 42, 123, 456 \\
Target language code & \texttt{nob\_Latn} \\
\bottomrule
\end{tabular}
\caption{Shared fine-tuning and generation configuration.}
\label{tab:training-config}
\end{table}

\section{Evaluation}

\subsection{Automatic MT metrics}

We report BLEU and chrF. BLEU is reported in the same fractional scale returned by the evaluation implementation; for example, 0.30 corresponds to approximately 30 BLEU in percentage-style reporting. chrF is reported on the usual 0--100 scale.

We do not report neural reference-based metrics such as COMET. This choice is deliberate: the goal of the paper is to analyze how fixed reference forms and written-standard-specific resources affect metric interpretation. BLEU, chrF, and exact-match terminology metrics are transparent enough to show when changes in score are driven by written-standard conformity rather than by translation adequacy alone.

We also report glossary-based terminology metrics. Term recall (TermR) checks whether English source terms with known Norwegian terminology equivalents are realized in the model output. To avoid treating recall as complete terminology accuracy, we additionally report term precision (TermP), which measures whether predicted glossary terms are licensed by source-side terms, and term F1 (TermF1), the harmonic mean of TermR and TermP. Term coverage (TermCov) is the proportion of source sentences containing at least one glossary term. The glossary contains 70 English--Norwegian term pairs. It covers 682 of 1,742 original-test sentences (39.15\%, 984 source-term instances) and 514 of 1,313 \nb{}-test sentences (39.15\%, 748 source-term instances). The glossary equivalents are \nb{} forms. Consequently, these metrics should be interpreted as \nb{}-oriented term-form conformity, not language-neutral terminology accuracy, especially on mixed-standard or \nn{}-reference data.

\subsection{Statistical testing}

For paired significance testing, we use 1,000-sample paired bootstrap resampling over sentence-level outputs. We apply this to sentence-level chrF and to the glossary-based terminology metrics. The main result tables still report corpus-level BLEU and chrF; the bootstrap test is used only to assess whether the direction of the observed differences is stable under sentence resampling. For SLIDE written-standard rates, we use exact McNemar tests on paired sentence-level labels, comparing whether each system output is NB-only or NN-only.

\subsection{SLIDE written-standard analysis}
\label{sec:slide-analysis}

To identify the written standard of model outputs, we apply SLIDE to both predictions and references. SLIDE is multi-label: a sentence may be compatible with \nb{}, \nn{}, both, or neither. We use a probability threshold of 0.5 and define:
\begin{itemize}
    \item NB-only: \texttt{nb} $\geq 0.5$ and \texttt{nn} $< 0.5$.
    \item NN-only: \texttt{nn} $\geq 0.5$ and \texttt{nb} $< 0.5$.
    \item Mixed: \texttt{nb} $\geq 0.5$ and \texttt{nn} $\geq 0.5$.
    \item Other/uncertain: neither \texttt{nb} nor \texttt{nn} reaches 0.5.
\end{itemize}

We report percentages for NB-only, NN-only, and mixed outputs, and the mean margin between \texttt{nb} and \texttt{nn} probabilities. This margin is useful because it captures how strongly a sentence is classified as \nb{} rather than \nn{}.

Because SLIDE is used both in target-side filtering and in output analysis, SLIDE percentages are classifier-based evidence rather than fully independent proof of written-standard shift. We therefore validate SLIDE on a stratified 120-item petroleum-domain sample drawn from references and model predictions. Table~\ref{tab:slide-validation} separates two questions: whether the manual annotation task is stable, and how closely SLIDE matches human labels.

\begin{table}[t]
\centering
\small
\resizebox{\columnwidth}{!}{%
\begin{tabular}{lccccc}
\toprule
Comparison & Exact & Macro-F1 & $\kappa$ & NB Prec. & NN Rec. \\
\midrule
Reviewer 1 vs Reviewer 2 & 84.2 & -- & 0.770 & -- & -- \\
Reviewer 1 vs SLIDE & 64.2 & 0.586 & 0.522 & 0.967 & 1.000 \\
Reviewer 2 vs SLIDE & 60.8 & 0.552 & 0.478 & 0.900 & 0.957 \\
\bottomrule
\end{tabular}%
}
\caption{Human validation of SLIDE written-standard labels on a stratified 120-item petroleum-domain sample. Exact agreement is reported as a percentage. NB Prec. is SLIDE precision for the NB-only label; NN Rec. is SLIDE recall for the NN-only label.}
\label{tab:slide-validation}
\end{table}

The two reviewers show substantial inter-annotator agreement, indicating that the manual labeling task is reasonably stable. Agreement between SLIDE and the human annotation passes is moderate rather than gold-standard level. Most disagreements involve the mixed and uncertain categories: SLIDE often assigns mixed labels to short, formulaic, or weakly diagnostic petroleum-domain sentences that human reviewers mark as uncertain or NB-only. For the two main contrastive labels used in our analysis, agreement is stronger: SLIDE NB-only precision is 0.900--0.967, and SLIDE NN-only recall is 0.957--1.000. We therefore interpret SLIDE as useful for detecting the large NB-only increase and NN-only suppression reported below, but not as a fine-grained gold standard for mixedness. The main MT conclusions do not rely on SLIDE alone: they are triangulated with BLEU/chrF reversals, \nb{}-oriented terminology metrics, and blind human judgments separating adequacy from \nb{} conformity.

\section{In-Domain Results}

\subsection{Size-Controlled In-Domain Effects}

Table~\ref{tab:indomain} shows the in-domain evaluation matrix. The apparent best system depends on the written standard of the reference set: the original model has the highest BLEU and chrF on the original mixed-standard test set, whereas the \nb{}-filtered model has the highest BLEU and chrF on the \nb{} test set. This is the first indication of evaluation bias: the same model behavior can be rewarded or penalized depending on which written standard the test set favors.

\begin{table*}[t]
\centering
\small
\begin{tabular}{llrrrrr}
\toprule
Train condition & Test set & BLEU & chrF & TermR & TermP & TermF1 \\
\midrule
Original & Original & 0.6290$_{\pm .0020}$ & 79.84$_{\pm .10}$ & 0.7497$_{\pm .0082}$ & 0.7030$_{\pm .0039}$ & 0.7256$_{\pm .0059}$ \\
Original & \nb{} & 0.6038$_{\pm .0016}$ & 78.70$_{\pm .06}$ & 0.8191$_{\pm .0074}$ & 0.7282$_{\pm .0047}$ & 0.7710$_{\pm .0058}$ \\
Original-subsampled & Original & 0.6177$_{\pm .0030}$ & 79.12$_{\pm .22}$ & 0.7605$_{\pm .0109}$ & 0.7035$_{\pm .0022}$ & 0.7309$_{\pm .0062}$ \\
Original-subsampled & \nb{} & 0.5937$_{\pm .0035}$ & 78.02$_{\pm .28}$ & 0.8200$_{\pm .0097}$ & 0.7261$_{\pm .0022}$ & 0.7702$_{\pm .0055}$ \\
\nb{} & Original & 0.5849$_{\pm .0005}$ & 77.43$_{\pm .11}$ & 0.8665$_{\pm .0016}$ & 0.7273$_{\pm .0011}$ & 0.7908$_{\pm .0013}$ \\
\nb{} & \nb{} & 0.6128$_{\pm .0001}$ & 79.34$_{\pm .12}$ & 0.8587$_{\pm .0015}$ & 0.7335$_{\pm .0008}$ & 0.7912$_{\pm .0011}$ \\
\bottomrule
\end{tabular}
\caption{In-domain evaluation. Values are mean$_{\pm}$standard deviation over three seeds. TermR, TermP, and TermF1 are glossary-based and \nb{}-oriented.}
\label{tab:indomain}
\end{table*}

The strongest test of the filtering effect compares \nb{} with original-subsampled, since both use 10,114 training examples. Under \nb{} references, filtering increases BLEU from 0.5937 to 0.6128 and chrF from 78.02 to 79.34. Under the original mixed-standard references, however, the same filtering lowers BLEU from 0.6177 to 0.5849 and chrF from 79.12 to 77.43. This reversal shows that the effect is not a uniform quality gain: \nb{} filtering improves \nb{}-conforming performance when the reference set rewards \nb{}, but harms string-overlap metrics when the reference distribution is mixed.

The terminology metrics show the complementary side of the same pattern. Because the glossary uses \nb{} term forms, higher terminology scores should be interpreted as increased \nb{} term-form conformity rather than language-neutral terminology accuracy. The \nb{}-filtered model increases TermF1 over original-subsampled on both the \nb{} test set (0.7702 to 0.7912) and the original test set (0.7309 to 0.7908). This indicates that filtering changes the written-standard form of domain terminology even when it hurts BLEU and chrF against mixed-standard references.

Paired bootstrap tests support the same interpretation: the \nb{} model significantly increases sentence-level chrF and \nb{}-oriented terminology scores on the \nb{} test set, but significantly decreases sentence-level chrF on the original mixed-standard test set while increasing \nb{} terminology conformity. Full test statistics are reported in Appendix~\ref{app:significance-tests}.

\subsection{Model-scale robustness}

Figure~\ref{fig:scale-trend} summarizes the size-controlled comparison across the 600M, 1.3B, and 3.3B NLLB-200 backbones. Two patterns remain consistent across all model sizes. First, \nb{} filtering improves BLEU on the \nb{} test set relative to the size-controlled original-subsampled baseline. Second, on the original mixed-standard test set, \nb{} filtering produces substantially more NB-only outputs. Although larger backbones improve absolute translation quality, they do not remove the evaluation reversal caused by target-standard mismatch. Full model-scale results, including chrF and TermF1, are reported in Appendix~\ref{app:model-scale-results}, Table~\ref{tab:scale}.

\begin{figure}[t]
\centering
\resizebox{\columnwidth}{!}{%
\begin{tikzpicture}
\begin{groupplot}[
  group style={group size=2 by 1, horizontal sep=1.6cm},
  width=4.2cm,
  height=3.8cm,
  xtick={1,2,3},
  xticklabels={600M,1.3B,3.3B},
  x tick label style={font=\scriptsize},
  y tick label style={font=\scriptsize},
  ylabel style={font=\scriptsize},
  title style={font=\scriptsize},
  legend to name=sharedlegend,
  legend columns=2,
  legend style={font=\scriptsize},
  every axis plot/.append style={thick, mark size=1.8pt},
]

\nextgroupplot[
  title={BLEU (\nb{} test set)},
  ylabel={BLEU},
  ymin=0.55, ymax=0.68,
]
\addplot[color=blue!60!black, mark=*] coordinates {(1,0.6128) (2,0.6313) (3,0.6447)};
\addlegendentry{\nb{}-filtered}
\addplot[color=gray!70!black, mark=square*, dashed] coordinates {(1,0.5937) (2,0.6078) (3,0.6271)};
\addlegendentry{Original-subsampled}

\nextgroupplot[
  title={NB-only outputs (\%), original test set},
  ylabel={NB-only \%},
  ymin=70, ymax=100,
]
\addplot[color=blue!60!black, mark=*] coordinates {(1,93.7) (2,94.3) (3,94.6)};
\addplot[color=gray!70!black, mark=square*, dashed] coordinates {(1,79.1) (2,79.6) (3,79.8)};

\end{groupplot}
\end{tikzpicture}%
}

\vspace{4pt}
\centering
\ref{sharedlegend}

\caption{Automatic-metric and written-standard shift patterns across NLLB-200 model sizes. \nb{}-filtered training consistently outperforms the size-controlled baseline on \nb{} BLEU and shows a stable elevated NB-only output rate on the original mixed-standard test set.}
\label{fig:scale-trend}
\end{figure}

The written-standard shift is equally robust across model scales. On the original mixed-standard test set, \nb{}-filtered models produce 14--15 percentage points more NB-only outputs than the size-controlled baseline at every scale (93.7--94.6\% vs.\ 79.1--79.8\%), while NN-only outputs are almost eliminated (13.7--14.2\% vs.\ 0.08--0.17\%). On the \nb{} test set, NB-only outputs further increase from 94.9--96.3\% to 99.0--99.4\%. Together, these results show that increasing model capacity improves absolute scores but does not reduce the tendency of \nb{}-filtered models to generate \nb{}, nor does it remove the evaluation bias introduced by written-standard mismatch.

\section{Written-Standard Shift}

Automatic MT metrics alone cannot determine whether the \nb{} model shifts the written-standard distribution of its outputs. Table~\ref{tab:slide} directly examines this question using SLIDE.

\begin{table*}[t]
\centering
\small
\begin{tabular}{llrrrr}
\toprule
Model / reference & Test & NB-only & NN-only & Mixed & Mean NB--NN \\
\midrule
Reference & Original & 76.2 & 17.2 & 6.0 & 0.583 \\
Original & Original & 79.0 & 14.1 & 5.9 & 0.643 \\
Original-subsampled & Original & 79.0 & 14.2 & 5.7 & 0.642 \\
\nb{} & Original & 93.4 & 0.7 & 4.8 & 0.919 \\
\midrule
Reference & \nb{} & 99.6 & 0.2 & 0.2 & 0.983 \\
Original & \nb{} & 95.1 & 3.6 & 1.0 & 0.905 \\
Original-subsampled & \nb{} & 94.8 & 4.1 & 0.7 & 0.897 \\
\nb{} & \nb{} & 98.8 & 0.03 & 0.8 & 0.976 \\
\bottomrule
\end{tabular}
\caption{SLIDE written-standard analysis of in-domain predictions and references. Percentages are averaged over seeds.}
\label{tab:slide}
\end{table*}

The original test references are mixed, containing 76.2\% \nb{}-only, 17.2\% \nn{}-only, and 6.0\% mixed sentences. The original and original-subsampled models approximately preserve this distribution, producing around 79\% \nb{}-only and 14\% \nn{}-only outputs. In contrast, the \nb{}-filtered model shifts strongly toward \nb{}, producing 93.4\% \nb{}-only and only 0.7\% \nn{}-only outputs on the same test set.

This provides the central evidence for RQ4. The \nb{} model does not simply improve or worsen automatic scores; it changes the written-standard preference expressed in its outputs. On the mixed original test set, it shifts toward \nb{} beyond the proportion observed in the references. This explains the earlier metric pattern: BLEU and chrF decrease because the references contain \nn{} and mixed-standard material, whereas glossary-based TermR and TermF1 increase because \nb{} term forms are rewarded.

On the \nb{} test set, the reference distribution is almost entirely \nb{}-only. The \nb{} model most closely matches this distribution, producing 98.8\% \nb{}-only outputs and essentially eliminating \nn{}-only outputs. This supports the interpretation that \nb{} filtering is effective when the desired deployment target is explicitly \nb{}, while also showing that the filtering procedure introduces a strong target-standard preference.

McNemar tests confirm that these SLIDE distribution shifts are statistically significant across seeds. Compared with original-subsampled, the \nb{} model increases \nb{}-only outputs on the original test set by 13.8--15.2 percentage points and reduces \nn{}-only outputs by 13.1--14.2 points. On the \nb{} test set, it further increases \nb{}-only outputs by 3.4--5.0 points and reduces \nn{}-only outputs by 3.5--4.9 points, with $p<0.001$ for every seed.

\section{Out-of-Domain Results}

Table~\ref{tab:flores} presents FLORES evaluation results. The zero-shot base model achieves higher BLEU and chrF scores than all fine-tuned systems on FLORES NB, indicating that domain fine-tuning reduces general-domain performance regardless of whether filtering is applied. Among the fine-tuned systems, differences are small: the \nb{} model is only marginally better than original-subsampled on FLORES NB and does not improve performance on FLORES NN.

\begin{table*}[t]
\centering
\small
\begin{tabular}{lrrrr}
\toprule
Model & FLORES NB BLEU & FLORES NB chrF & FLORES NN BLEU & FLORES NN chrF \\
\midrule
Base NLLB & 0.3059 & 59.22 & 0.1713 & 50.52 \\
Original & 0.2726$_{\pm .0015}$ & 56.35$_{\pm .23}$ & 0.1524$_{\pm .0006}$ & 48.28$_{\pm .15}$ \\
Original-subsampled & 0.2750$_{\pm .0009}$ & 56.56$_{\pm .01}$ & 0.1543$_{\pm .0006}$ & 48.48$_{\pm .04}$ \\
\nb{} & 0.2766$_{\pm .0025}$ & 56.63$_{\pm .22}$ & 0.1537$_{\pm .0007}$ & 48.33$_{\pm .11}$ \\
\bottomrule
\end{tabular}
\caption{Out-of-domain FLORES evaluation. Fine-tuned systems show mean$_{\pm}$standard deviation over three seeds; the base model is run once. All systems decode with \texttt{nob\_Latn}.}
\label{tab:flores}
\end{table*}

The FLORES results provide a limited answer to RQ5. The \nb{}-conforming effects observed in-domain do not substantially transfer to the general-domain FLORES benchmark. Instead, the main finding concerns written-standard bias in evaluation. FLORES NB references are 99.5\% \nb{}-only, while FLORES NN references are 99.5\% \nn{}-only. However, all fine-tuned systems generate approximately 97--98\% \nb{}-only outputs on both FLORES NB and FLORES NN. This is expected because all systems decode with \texttt{nob\_Latn}, which is \nb{}-oriented. Therefore, FLORES NN scores should not be interpreted as measuring \nn{} translation ability; rather, they measure the metric penalty incurred when \nb{}-oriented outputs are evaluated against \nn{} references.

\section{Discussion}

The results support the view that \nb{} filtering induces target-standard bias toward the dominant written standard. As defined in Section~\ref{sec:bias}, this bias is not the same as harmful specialization in every setting. For many industrial deployments, \nb{} is the correct and expected target. The interpretation depends on the deployment goal: if the desired output is explicitly \nb{} petroleum translation, filtering is useful because it improves \nb{}-conforming in-domain performance relative to a size-controlled baseline. If the goal is general Norwegian MT or robustness to a mixed-standard reference distribution, the same filtering is harmful because it lowers BLEU and chrF on the original test set and suppresses \nn{}-only outputs.

\paragraph{Evaluation bias and dominant-standard preference.}
The experiments also expose an evaluation bias that is specific to multi-standard languages. Reference-based MT metrics do not distinguish adequacy errors from legitimate written-standard differences. A \nb{} output can receive lower BLEU or chrF against a \nn{}-like reference even when it is semantically adequate, while a \nb{}-oriented glossary can make the same output look better in terminology metrics. This means that evaluation design implicitly encodes a dominant-standard preference. The effect is visible in both directions in our results: \nb{} filtering lowers string-overlap scores on the mixed original test set, but increases \nb{}-oriented terminology scores and SLIDE \nb{} conformity. In a \nb{} deployment scenario, TermR, TermP, and TermF1 remain useful diagnostics because they measure whether the system uses expected \nb{} terminology forms and whether those forms are source-licensed; but in a mixed-standard scenario, the same metrics can overstate general quality by rewarding \nb{} forms even when the reference is \nn{} or mixed. This is not a flaw of the metric if its purpose is clearly stated, but a risk if presented as language-neutral terminology accuracy. We therefore treat written-standard identification as part of evaluation, not only as a preprocessing diagnostic.

\paragraph{Industrial relevance.}
The industrial relevance of target-standard control is visible in commercial MT interfaces and APIs. Several major MT providers expose Norwegian translation through \nb{}-oriented labels or language codes: DeepL lists Norwegian as \nb{}-only and uses \texttt{NB} in its API \cite{deepl-languages,deepl-help-languages}; Microsoft Translator lists Norwegian \nb{} with code \texttt{nb} \cite{microsoft-translator-languages}; and Google Cloud Translation exposes a Norwegian Bokmal code \texttt{nb}, alongside a more general Norwegian code \cite{google-translate-languages}. Petroleum-domain terminology resources are also \nb{}-oriented in practice; for example, the Petroleum Safety Authority Norway's terminology list gives Norwegian industry terms rather than parallel \nb{}/\nn{} term variants \cite{havtil-terms}. This supports our interpretation of TermR, TermP, and TermF1 as \nb{} term-form conformity metrics for a \nb{} deployment setting. The deployment choice itself is legitimate. The audit problem is whether systems described as ``Norwegian'' are actually general Norwegian, \nb{}-specialized, or \nn{}-capable.

The FLORES results add another caution. Fine-tuning specializes the 600M model for the petroleum domain, but it does not improve general-domain performance. The base NLLB model is stronger on FLORES NB than all fine-tuned variants. Thus, the paper should not claim that \nb{} filtering improves English--Norwegian MT in general. The supported claim is narrower and more important for MT audit: filtering changes what the model considers normal Norwegian, while changing robustness and metric interpretation. The larger-model replication strengthens the in-domain part of this claim by showing that both the BLEU/chrF reversal and the written-standard shift are not artifacts of the 600M backbone, while the out-of-domain diagnostics remain reported for the full 600M matrix.

More broadly, the study shows that generated-language systems can encode normative preferences through seemingly technical design choices. In MT, these preferences appear as written-standard specialization; in other generation settings, analogous effects may appear as framing, emphasis, or omission.

\section{Human Evaluation}

To validate the automatic written-standard analysis, we conducted a blind diagnostic human evaluation on 100 sentence-level items: 50 written-standard shift cases and 50 control cases. Because shift cases were intentionally enriched, the evaluation tests whether the detected shifts are human-perceptible; it does not estimate population-level preference rates. Two annotators rated both outputs for adequacy and \nb{} conformity on a 0--2 scale and selected the preferred output for a \nb{} petroleum-domain translation use case. System identities were hidden during annotation. One item in the second annotation pass was invalid, leaving 99 shared valid items for inter-annotator agreement.

Table~\ref{tab:human-eval} summarizes the results. Across both annotation passes, adequacy differences between the \nb{}-filtered and original-subsampled systems are small (+0.01--+0.02), while \nb{} conformity increases substantially (+0.79--+0.98). Preference judgments also favor the \nb{}-filtered system, but primarily in the shift-enriched subset; control cases contain many ties and show no consistent advantage for the filtered system. This supports the interpretation that filtering mainly changes target-standard conformity rather than semantic adequacy.

\begin{table}[t]
\centering
\small
\resizebox{\columnwidth}{!}{%
\begin{tabular}{lrrrrr}
\toprule
Pass & $n$ & $\Delta$Adeq. & $\Delta$\nb{} & \nb{} pref. & Orig. pref. \\
\midrule
Annotator 1 & 100 & +0.02 & +0.98 & 59 & 16 \\
Annotator 2 & 99 & +0.01 & +0.79 & 55 & 11 \\
\bottomrule
\end{tabular}%
}
\caption{Blind diagnostic human evaluation comparing the \nb{}-filtered model with the size-controlled original-subsampled model. $\Delta$ values are \nb{} minus original-subsampled mean scores on 0--2 scales. Preference counts exclude ties; there are 25 ties for Annotator 1 and 33 ties for Annotator 2.}
\label{tab:human-eval}
\end{table}

Inter-annotator agreement on the 99 shared valid items is substantial for the overall \nb{}-deployment preference (85.9\% exact agreement, $\kappa=0.751$). Agreement is also substantial for \nb{} conformity judgments and moderate to substantial for adequacy judgments. At the item level, the automatic SLIDE shift margin correlates strongly with the human-rated \nb{} conformity gain (Spearman's $\rho=0.762$) and moderately with human \nb{}-deployment preference ($\rho=0.676$). The human evaluation therefore triangulates the automatic analysis: the selected shift cases correspond to human-perceived written-standard differences, not merely classifier artifacts.

Representative cases are provided in Appendix~\ref{app:qualitative-examples}.

\section{Ethical Considerations: Linguistic Inclusion}

This study is not an ethical evaluation of Norwegian MT systems in deployment, but the results have ethical implications for linguistic inclusion. \nb{} and \nn{} are both official written standards of Norwegian, yet they are not equally represented in many digital corpora and MT pipelines. If a system is trained, decoded, and evaluated under a generic ``Norwegian'' label, it may silently privilege \nb{} and make \nn{}-like outputs appear erroneous even when they are valid Norwegian.

The risk is not limited to training data. Evaluation choices also encode language-standard preferences. A \nb{} reference set, a \nb{}-oriented terminology glossary, or a forced \texttt{nob\_Latn} decoding setting can all make \nb{}-conforming outputs look objectively better. Conversely, evaluating a \nb{}-oriented system against \nn{} references can make valid \nb{} outputs look worse. In both cases, the metric result partly reflects written-standard bias rather than translation adequacy alone.

For deployment, the ethical requirement is transparency rather than neutrality. We do not claim that \nb{} specialization is inherently harmful; it is appropriate when the intended target is explicitly \nb{}. The problem is silent specialization. The practical issue is auditability: users, clients, and evaluators should be able to know whether a system is general Norwegian, \nb{}-specialized, or \nn{}-capable. A model trained for \nb{} should be described as \nb{}-specialized, not as general Norwegian MT. \nn{}-oriented or mixed-standard use cases require their own data, references, and terminology resources.

Turning this principle into practice does not require a single universal fix, but several concrete safeguards follow directly from our results. Training data can be balanced across \nb{} and \nn{} when the intended system is general Norwegian, or explicitly specialized when the intended output is \nb{}. Target-standard tags can make generation controllable rather than silent. Evaluation can report written-standard distributions alongside BLEU and chrF, and use standard-specific or dual-standard references where available. Terminology evaluation should also separate semantic term coverage from written-standard form conformity: if stable \nn{} terminology variants exist for a deployment setting, they can be added to a dual-standard glossary; if the available industry terminology resource is \nb{}-oriented, the metric should instead be explicitly reported as \nb{} term-form conformity. These steps would turn hidden target-standard preference into an auditable property of the MT pipeline, and are a practical step toward more transparent and inclusive MT development for languages with multiple legitimate written standards.

\section{Limitations}

This study has several limitations. First, it studies one translation direction and one domain. Although the empirical study is limited to Norwegian, the mechanism is not Norwegian-specific: target-side filtering, forced decoding codes, standard-specific references, and terminology glossaries are common MT design choices. We therefore treat Norwegian as a controlled case study of a broader evaluation problem in multi-standard languages, not as evidence that the same magnitude of bias will occur for every language pair. Second, all systems decode with \texttt{nob\_Latn}; therefore, FLORES NN is an evaluation-bias test rather than an evaluation of \nn{}-targeted MT. Third, the model-scale replication covers in-domain automatic metrics and SLIDE written-standard analysis; FLORES evaluation, significance testing, and human evaluation are reported for the full 600M diagnostic matrix. Fourth, the terminology metrics are based on \nb{} terminology forms and should be interpreted accordingly. Although we report precision and F1 in addition to recall, the metrics are still automatic exact-match diagnostics rather than substitutes for human terminology assessment. Fifth, the \nb{} filtering step may remove examples for reasons correlated with but not identical to written standard, such as style, sentence length, or domain subtopic. Sixth, the threshold sensitivity analysis checks retained data size, but does not retrain models under every threshold. Seventh, human validation is limited in scope: our blind evaluation covers a small diagnostic sample and is intended to check adequacy and written-standard conformity rather than estimate full-corpus preference rates or replace full human MT evaluation. Finally, the paired bootstrap test uses sentence-level chrF rather than corpus chrF; corpus-level chrF remains the primary reported score in the main tables.

\section{Conclusion}

This paper shows that target-side \nb{} filtering for English--Norwegian domain MT is not merely data cleaning. It can induce target-standard bias toward the dominant written standard. Under a data-size controlled comparison, \nb{} filtering increases in-domain \nb{}-conforming performance, but it reduces BLEU and chrF on the original mixed-standard distribution; both this automatic metric pattern and the SLIDE-measured written-standard shift persist across 600M, 1.3B, and 3.3B NLLB-200 backbones. Out-of-domain FLORES evaluation for the 600M systems shows that these effects are domain-bound and that \nn{} reference evaluation is dominated by written-standard bias. For languages with multiple written standards, MT data curation and evaluation should therefore make target-standard assumptions explicit and auditable.

\section*{Data and Code Availability}

We will release the data-processing scripts, fine-tuning and evaluation configurations, processed train/dev/test split metadata, SLIDE written-standard annotations, metric summaries, significance-test outputs, and anonymized human-evaluation materials. Large model checkpoints are not redistributed; instead, the release will provide the scripts and configuration needed to reproduce the fine-tuning runs from the public NLLB-200 checkpoints.

% \section*{Acknowledgements}
% Add HPC/funding acknowledgements here.

\bibliography{eamt26}
\bibliographystyle{eamt26}

\appendix

\section{Qualitative Examples}
\label{app:qualitative-examples}

Table~\ref{tab:examples} shows three representative cases from the blind human evaluation. The examples illustrate a written-standard shift, a control case, and a trade-off between \nb{} conformity and terminology precision.

\begin{table*}[h]
\centering
\small
\begin{tabular}{p{0.18\linewidth}p{0.24\linewidth}p{0.24\linewidth}p{0.26\linewidth}}
\toprule
Case & Original-subsampled & \nb{}-filtered & Human judgment \\
\midrule
Shift: production closure
& Produksjonen frå Vale og Volund var også stengd på grunn av tekniske problem.
& Produksjonen fra Vale og Volund ble også stengt på grunn av tekniske problemer.
& Both are adequate, but the baseline keeps \nn{}-like forms (\emph{frå}, \emph{stengd}); the filtered output is consistent \nb{}. \\
\midrule
Control: CO2 storage atlas
& Atlaset gir en oversikt over områder der CO2 kan lagres trygt i undergrunnen i lang tid.
& Atlaset gir oversikt over områder der CO2 kan lagres trygt i undergrunnen i lang tid.
& Both are adequate \nb{} outputs; annotators judged this as a tie. \\
\midrule
Trade-off: barrier weakness
& I juli var produksjonen frå Valefeltet stengd på grunn av tekniske problem og Gimle- og Tordisfeltet stengd på grunn av barriere svakhet i brønn.
& I juli ble produksjonen fra Vale-feltet stengt på grunn av tekniske problemer og feltene Gimle og Tordis stengt på grunn av barrierefølsomhet i brønn.
& The filtered output is more consistent \nb{}, but \emph{barrierefølsomhet} is less precise for ``barrier weakness'' than \emph{svakhet}. \\
\bottomrule
\end{tabular}
\caption{Representative cases from the blind human evaluation.}
\label{tab:examples}
\end{table*}

\section{Model-Scale Results}
\label{app:model-scale-results}

Table~\ref{tab:scale} reports the full size-controlled in-domain automatic metric replication across NLLB-200 backbones.

\begin{table*}[h]
\centering
\small
\begin{tabular}{lllrrr}
\toprule
Size & Test set & Model & BLEU & chrF & TermF1 \\
\midrule
\multirow{4}{*}{600M}
& \multirow{2}{*}{\nb{}} & Original-sub. & 0.5937 & 78.02 & 0.7702 \\
& & \nb{} & 0.6128 & 79.34 & 0.7912 \\
& \multirow{2}{*}{Original} & Original-sub. & 0.6177 & 79.12 & 0.7309 \\
& & \nb{} & 0.5849 & 77.43 & 0.7908 \\
\midrule
\multirow{4}{*}{1.3B}
& \multirow{2}{*}{\nb{}} & Original-sub. & 0.6078 & 79.08 & 0.7686 \\
& & \nb{} & 0.6313 & 80.50 & 0.7899 \\
& \multirow{2}{*}{Original} & Original-sub. & 0.6325 & 80.17 & 0.7286 \\
& & \nb{} & 0.5999 & 78.46 & 0.7905 \\
\midrule
\multirow{4}{*}{3.3B}
& \multirow{2}{*}{\nb{}} & Original-sub. & 0.6271 & 80.12 & 0.7759 \\
& & \nb{} & 0.6447 & 81.23 & 0.7856 \\
& \multirow{2}{*}{Original} & Original-sub. & 0.6515 & 81.21 & 0.7302 \\
& & \nb{} & 0.6091 & 79.03 & 0.7872 \\
\bottomrule
\end{tabular}
\caption{Size-controlled in-domain automatic metric replication across NLLB-200 backbones. Values are means over three seeds. Original-sub. denotes original-subsampled. Full standard deviations are reported in the released result summaries.}
\label{tab:scale}
\end{table*}

\section{Threshold Sensitivity Analysis}
\label{app:threshold-sensitivity}

To examine the robustness of the filtering procedure, we recomputed retained training-set sizes under nearby SLIDE threshold settings without retraining MT models. Table~\ref{tab:threshold-sensitivity} shows the resulting data-selection grid. Across \texttt{nb} thresholds from 0.60 to 0.90 and \texttt{nn} exclusion thresholds from 0.10 to 0.40, retained data ranges from 9,925 to 10,212 training sentences, corresponding to 71.2--73.3\% of the original training split. This indicates that the selected 0.80/0.30 threshold lies in a relatively stable region of the SLIDE score distribution. We treat this as a data-selection robustness check only; it does not show that model behavior would be unchanged under every threshold setting.

\begin{table}[h]
\centering
\small
\begin{tabular}{ccrr}
\toprule
\texttt{nb} threshold & \texttt{nn} threshold & Retained & Retained \% \\
\midrule
0.60 & 0.10 & 9,947 & 71.38 \\
0.60 & 0.20 & 10,053 & 72.14 \\
0.60 & 0.30 & 10,127 & 72.67 \\
0.60 & 0.40 & 10,212 & 73.28 \\
0.70 & 0.10 & 9,945 & 71.37 \\
0.70 & 0.20 & 10,049 & 72.11 \\
0.70 & 0.30 & 10,121 & 72.63 \\
0.70 & 0.40 & 10,203 & 73.22 \\
0.80 & 0.10 & 9,940 & 71.33 \\
0.80 & 0.20 & 10,041 & 72.06 \\
0.80 & 0.30 & 10,113 & 72.57 \\
0.80 & 0.40 & 10,193 & 73.15 \\
0.90 & 0.10 & 9,925 & 71.22 \\
0.90 & 0.20 & 10,023 & 71.93 \\
0.90 & 0.30 & 10,087 & 72.39 \\
0.90 & 0.40 & 10,157 & 72.89 \\
\bottomrule
\end{tabular}
\caption{Sensitivity of retained training-set size to SLIDE filtering thresholds. This check recomputes selection counts from the scored original training set and does not retrain MT models.}
\label{tab:threshold-sensitivity}
\end{table}

\section{Significance Tests}
\label{app:significance-tests}

Table~\ref{tab:significance-tests} reports paired bootstrap tests comparing the \nb{}-filtered model against the original-subsampled baseline. We use 1,000 resamples over paired sentence-level outputs. Positive deltas favor the \nb{}-filtered model.

\begin{table*}[h]
\centering
\small
\begin{tabular}{llrrr}
\toprule
Test set & Metric & Delta & 95\% CI & $p$ \\
\midrule
\nb{} & sentence chrF & +1.0585 & [+0.8561, +1.2522] & $<0.001$ \\
\nb{} & TermR & +0.0388 & [+0.0288, +0.0497] & $<0.001$ \\
\nb{} & TermP & +0.0074 & [+0.0021, +0.0129] & 0.006 \\
\nb{} & TermF1 & +0.0210 & [+0.0147, +0.0282] & $<0.001$ \\
Original & sentence chrF & -1.7956 & [-2.0872, -1.5148] & $<0.001$ \\
Original & TermR & +0.1060 & [+0.0920, +0.1197] & $<0.001$ \\
Original & TermP & +0.0238 & [+0.0182, +0.0298] & $<0.001$ \\
Original & TermF1 & +0.0599 & [+0.0517, +0.0685] & $<0.001$ \\
\bottomrule
\end{tabular}
\caption{Paired bootstrap tests for the size-controlled comparison between the \nb{}-filtered model and original-subsampled. Tests use sentence-level chrF and glossary-based terminology metrics.}
\label{tab:significance-tests}
\end{table*}

\end{document}
